% SPDX-License-Identifier: CC-BY-NC-ND-4.0
% !TEX root = main.tex

\section{加法圏}

\begin{definition} \label{def: IsPreadditive}
局所小圏$\mathscr{C}$は次の条件を満たすとき\term{ぜんかほうけん}{前加法圏}{preadditive category}であるという.
\begin{itemize}
    \item 任意の$A, B \in \Ob(\mathscr{C})$に対して写像$\sigma_{(A,B)} \colon \Mor_{\mathscr{C}}(A, B) \times \Mor_{\mathscr{C}}(A, B) \to \Mor_{\mathscr{C}}(A, B)$が定まり次の条件を満たす.\
    ただし任意の$f, g \in \Mor_{\mathscr{C}}(A, B)$に対して$f + g \coloneq \sigma_{(A,B)}(f,g)$とする.
    \begin{itemize}
    \item ある$0_{A,B}^{+} \in \Mor_{\mathscr{C}}(A,B)$が存在して$(\Mor_{\mathscr{C}}(A,B), \sigma_{(A,B)}, 0_{A,B} )$はアーベル群.
    \item 任意の$A, B, C \in \Ob(\mathscr{C})$と$f,g \in \Mor_{\mathscr{C}}(A, B),\ h \in \Mor_{\mathscr{C}}(B, C)$に対して$h \circ (f + g) = h \circ f + f \circ g$成り立つ.
    \item 任意の$A, B, C \in \Ob(\mathscr{C})$と$f,g \in \Mor_{\mathscr{C}}(B, C),\ h \in \Mor_{\mathscr{C}}(A, B)$に対して$(f + g) \circ h = f \circ h + g \circ h$成り立つ.
    \end{itemize}
\end{itemize}
\end{definition}

\begin{lemma} \label{lem: additive_zero_comp}
$\mathscr{C}$を零対象$Z$を持つ前加法圏とする.\
任意の$A, B, C \in \Ob(\mathscr{C})$と$g \in \Mor_{\mathscr{C}}(B, C),\ h \in \Mor_{\mathscr{C}}(C, A)$に対して次の図式が可換になる.
% https://q.uiver.app/#q=WzAsNixbMCwwLCJCIl0sWzEsMCwiQSJdLFsxLDEsIkMiXSxbMywwLCJCIl0sWzQsMCwiQSJdLFs0LDEsIkMiXSxbMSwwLCIwX3tBLEJ9XnsrfSIsMl0sWzAsMiwiZyIsMl0sWzEsMiwiMF97QSxDfV57K30iXSxbNSwzLCIwX3tDLEJ9XnsrfSJdLFs0LDMsIjBfe0EsIEJ9XnsrfSIsMl0sWzUsNCwiaCIsMl1d
\[\begin{tikzcd}
	B & A && B & A \\
	& C &&& C
	\arrow["g"', from=1-1, to=2-2]
	\arrow["{0_{A,B}^{+}}"', from=1-2, to=1-1]
	\arrow["{0_{A,C}^{+}}", from=1-2, to=2-2]
	\arrow["{0_{A, B}^{+}}"', from=1-5, to=1-4]
	\arrow["{0_{C,B}^{+}}", from=2-5, to=1-4]
	\arrow["h"', from=2-5, to=1-5]
\end{tikzcd}\]
\end{lemma}

\begin{proof}
$g \circ 0_{A,B}^{+} = g \circ (0_{A,B}^{+} + 0_{A,B}^{+}) = g \circ 0_{A,B}^{+} + g \circ 0_{A,B}^{+}$であり両辺$g \circ 0_{A,B}^{+}$を引けば$0_{A,C}^{+} = g \circ 0_{A,B}^{+}$となる.\
同様に$0_{A,B}^{+} \circ h = 0_{C,B}^{+}$となる.
\end{proof}

\begin{proposition} \label{prop: additive_zero_eq_zero_morphism}
$\mathscr{C}$を零対象$Z$を持つ前加法圏とする.\
任意の$A,B \in \Ob(\mathscr{C})$に対して次が成り立つ.
\begin{center}
    $0_{A,B}^{+} = 0_{A,B}$
\end{center}
\end{proposition}

\begin{proof}
$\Mor_{\mathscr{C}}(Z, B)$はシングルトンなので$\iexcl_B = 0_{Z,B}^{+}$である.\
\cref{lem: additive_zero_comp}により$0_{A,B} = !_A \circ 0_{Z,B}^{+} = 0_{A,B}^{+}$となる.
\end{proof}

\begin{definition} \label{def: IsAdditive}
前加法圏$\mathscr{C}$は次の条件を満たすとき\term{かほうけん}{加法圏}{additive category}であるという.
\begin{itemize}
    \item 零対象$Z \in \Ob(\mathscr{C})$が存在.
    \item 任意の$A, B \in \Ob(\mathscr{C})$に対して$D_{A,B} \colon \mathbf{2} \to \mathscr{C}$の極限$(A \times B, \pi)$が存在.
\end{itemize}
\end{definition}

\begin{proposition} \label{prop: direct_product_and_direct_sum_of_biproduct}
$\mathscr{C}$を前加法圏とし$A, B \in \Ob(\mathscr{C})$とする.\
$(A \oplus B, \pi, \iota)$が\cref{def: IsBiproduct}の図式を満たし$\sum_{i \in 2} \iota_i \circ \pi_i = \id_{A \oplus B}$が成り立つことと$(A \oplus B, \pi, \iota)$が双積であることは同値である.\
\end{proposition}

\begin{proof}
$X \in \Ob(\mathscr{C})$とし$f \in \Mor_{\mathscr{C}}(X, A),\ g \in \Mor_{\mathscr{C}}(X, B)$とする.\
$h \coloneq \iota_0 \circ f + \iota_1 \circ g$とすると
\begin{align*}
\pi_0 \circ h &= \pi_0 \circ \iota_0 \circ f + \pi_0 \circ \iota_1 \circ g \\
	      &= \id_{A} \circ f + 0_{B, A} \circ g \\
	      &= f + 0_{X, A} \\
	      &= f
\end{align*}
となり同様に$\pi_1 \circ h = g$となる.

$\pi_0 \circ h' = f$かつ$\pi_1 \circ h' = g$とすると
\begin{align*}
h' &= \id_{A \oplus B} \circ h' \\
   &= (\iota_0 \circ  \pi_0 + \iota_1 \circ \pi_1) \circ h' \\
   &= \iota_0 \circ f + \iota_1 \circ g \\
   &= h
\end{align*}
となるので$(A \oplus B, \pi)$は直積である.\

$X \in \Ob(\mathscr{C})$とし$f \in \Mor_{\mathscr{C}}(A, X),\ g \in \Mor_{\mathscr{C}}(B, X)$とする.\
$h \coloneq f \circ \pi_0 + g \circ \pi_1$とすると
\begin{align*}
h \circ \iota_0 &= f \circ \pi_0 \circ \iota_0 + g \circ \pi_1 \circ \iota_0 \\
		&= f \circ \id_{A} + g \circ 0_{A, B} \\
		&= f + 0_{A, X} \\
		&= f
\end{align*}
となり同様に$h \circ \iota_1 = g$となる.

$h' \circ \iota_0 = f$かつ$h' \circ \iota_1 = g$とすると
\begin{align*}
h' &= h' \circ \id_{A \oplus B} \\
   &= h' \circ (\iota_0 \circ \pi_0 + \iota_1 \circ \pi_1) \\
   &= f \circ \pi_0 + g \circ \pi_1 \\
   &= h
\end{align*}
となるので$(A \oplus B, \iota)$は直和である.\

$(A \oplus B, \pi)$が直積で$(A \oplus B, \iota)$が直和のとき$e \coloneq \iota_0 \circ \pi_0 + \iota_1 \circ \pi_1$とすると$i \in 2$に対して$\pi_i \circ e = \pi_i$が成り立つので極限の一意性により$e = \id_{A \oplus B}$である.
\end{proof}

\begin{proposition} \label{prop: biproduct_of_direct_product}
$\mathscr{C}$を前加法圏とし$A, B \in \Ob(\mathscr{C})$とする.\
$(A \times B, \pi)$を直積とする.\
このとき$\iota_0 \coloneq \langle \id_A, 0_{A, B} \rangle$とし$\iota_1 \coloneq \langle 0_{B, A}, \id_B \rangle$とすると$(A \times B, \pi, \iota)$は双積となる.
\end{proposition}

\begin{proof}
$\langle -,- \rangle$の定義(\cref{def: directProduct})から次の図式が可換になる.
% https://q.uiver.app/#q=WzAsNSxbMCwxLCJCIl0sWzIsMCwiQSJdLFsyLDEsIkEgXFx0aW1lcyBCIl0sWzIsMiwiQiJdLFs0LDEsIkEiXSxbNCwyLCJcXGxhbmdsZSBcXG1hdGhybXtpZH1fQSwgMF97QSwgQn0gXFxyYW5nbGUiLDJdLFsyLDEsIlxccGlfMCJdLFs0LDEsIlxcbWF0aHJte2lkfV9BIiwyLHsiY3VydmUiOjJ9XSxbMCwyLCJcXGxhbmdsZSAwX3tCLCBBfSwgXFxtYXRocm17aWR9X0IgXFxyYW5nbGUiLDJdLFsyLDMsIlxccGlfMSJdLFswLDMsIlxcbWF0aHJte2lkfV9CIiwyLHsiY3VydmUiOjJ9XSxbMCwxLCIwX3tCLEF9IiwwLHsiY3VydmUiOi0yfV0sWzQsMywiMF97QSxCfSIsMCx7ImN1cnZlIjotMn1dXQ==
\[\begin{tikzcd}
	&& A && \\
	B && {A \times B} && A \\
	&& B
	\arrow["{0_{B,A}}", curve={height=-12pt}, from=2-1, to=1-3]
	\arrow["{\langle 0_{B, A}, \mathrm{id}_B \rangle}"', from=2-1, to=2-3]
	\arrow["{\mathrm{id}_B}"', curve={height=12pt}, from=2-1, to=3-3]
	\arrow["{\pi_0}", from=2-3, to=1-3]
	\arrow["{\pi_1}", from=2-3, to=3-3]
	\arrow["{\mathrm{id}_A}"', curve={height=12pt}, from=2-5, to=1-3]
	\arrow["{\langle \mathrm{id}_A, 0_{A, B} \rangle}"', from=2-5, to=2-3]
	\arrow["{0_{A,B}}", curve={height=-12pt}, from=2-5, to=3-3]
\end{tikzcd}\]

$e \coloneq \iota_0 \circ \pi_0 + \iota_1 \circ \pi_1$とする.\
$\pi_0 \circ e = \pi_0 \circ \iota_0 \circ \pi_0 + \pi_0 \circ \iota_1 \circ \pi_1 = \pi_0$であり同様に$\pi_1 \circ e = \pi_1$となる.\
$\pi_0 \circ \id_{A \times B} = \pi_0$かつ$\pi_1 \circ \id_{A \times B} = \pi_1$なので一意性から$e = \id_{A \times B}$である.
\end{proof}

\begin{proposition} \label{prop: biproduct_of_direct_sum}
$\mathscr{C}$を前加法圏とし$A, B \in \Ob(\mathscr{C})$とする.\
$(A \amalg B, \iota)$を直和する.\
このとき$\pi_0 \coloneq \lbrack \id_A, 0_{A, B} \rbrack$とし$\pi_1 \coloneq \lbrack 0_{B, A}, \id_B \rbrack$とすると$(A \amalg B, \pi, \iota)$は双積となる.
\end{proposition}

\begin{proof}
\cref{prop: biproduct_of_direct_product}の双対である.
\end{proof}
